\section{The receding-horizon idea}
\label{sec:intro}

\subsection{What a controller has to do}

A \emph{process} is a system whose evolution in time we want to influence: a tank whose
level has to be kept constant, a storage system whose charge has to follow the needs of a
building, a traffic light whose colours have to keep the queues short. We act on the process
through its \emph{control inputs} and we observe it through its \emph{measurements}. A
\emph{controller} is the rule that, at every sampling time, turns the measurements into the next
value of the inputs.

Classical controllers, such as the PID, compute the input from the current error with a fixed
formula. They are simple and robust, but they do not know the limits of the process --- a pump
that cannot exceed its rated power, a tank that cannot be emptied below a minimum level --- and
they do not look ahead. Model predictive control (MPC) is a different way of writing a controller:
at every sampling time it \emph{plans} the inputs over a window of future steps by solving an
optimization problem, and it uses only the first step of the plan.

\subsection{The receding horizon}

Fix a sampling time $T$ and call $k = 0, 1, 2, \dots$ the sampling instants. At time $k$ the
controller measures the state $x_k$ of the process, and considers the next $N$ steps, the
\emph{prediction horizon}. Inside the window it uses a second time index, $i = 0, 1, \dots, N$,
relative to the beginning of the window: the quantity predicted for step $i$ of the window opened
at time $k$ is written with the double index $i|k$, as in $x_{\ik{i}}$ and $u_{\ik{i}}$. The
absolute time of that step is $k + i$.

The controller then solves an optimization problem whose unknowns are the inputs over the
window, $u_{\ik{0}}, \dots, u_{\ik{N-1}}$, and whose cost measures how well the predicted
trajectory meets the control goal. Of the optimal sequence only the first element is applied to
the process, $u_k = u^*_{\ik{0}}$. At time $k+1$ the state is measured again and the whole
procedure is repeated over a window that has moved one step forward. This is the
\emph{receding horizon}: the window recedes as time advances (\cref{fig:receding}).

\begin{figure}[t]
\centering
\begin{tikzpicture}[x=0.9cm,y=0.9cm,>=Stealth]
  % time axis
  \draw[->] (-0.5,0) -- (13.5,0) node[right] {time};
  \foreach \t in {0,...,10} \draw (\t,0.08) -- (\t,-0.08);
  \node[below] at (3,-0.1) {$k$};
  \node[below] at (4,-0.1) {$k+1$};
  \node[below] at (8,-0.1) {$k+N$};
  % window at k
  \draw[very thick,mpc] (3,1.5) -- (8,1.5);
  \draw[mpc] (3,1.3) -- (3,1.7);
  \draw[mpc] (8,1.3) -- (8,1.7);
  \node[right,mpc,font=\small] at (8.2,1.5) {window opened at $k$: $i = 0,\dots,N$};
  \fill[mpc] (3,1.5) circle (0.09);
  \node[left,mpc,font=\small] at (2.8,1.5) {applied: $u^*_{0|k}$};
  % window at k+1
  \draw[very thick,sap] (4,0.75) -- (9,0.75);
  \draw[sap] (4,0.55) -- (4,0.95);
  \draw[sap] (9,0.55) -- (9,0.95);
  \fill[sap] (4,0.75) circle (0.09);
  \node[right,sap,font=\small] at (9.2,0.75) {window opened at $k+1$};
  % predicted trajectory sketch
  \draw[thick,mpc,dashed] plot[smooth] coordinates {(3,3.0) (4,2.7) (5,2.45) (6,2.35) (7,2.3) (8,2.3)};
  \draw[thick] plot[smooth] coordinates {(0,3.6) (1,3.4) (2,3.25) (3,3.0)};
  \fill (3,3.0) circle (0.07);
  \node[above,font=\small] at (1.2,3.45) {measured past};
  \node[right,mpc,font=\small] at (8.1,2.3) {predicted};
  \node[left,font=\small] at (-0.1,3.0) {$x$};
\end{tikzpicture}
\caption{The receding horizon. At time $k$ the controller measures $x_k$, plans the inputs over
the next $N$ steps, applies the first one and discards the rest. At $k+1$ it starts again with a
window moved one step ahead.}
\label{fig:receding}
\end{figure}

\subsection{The nominal MPC algorithm}

The procedure can be written as an algorithm. The adjective \emph{nominal} means that the
controller trusts its model completely: no disturbances, no model errors, a state that can be
measured.

\begin{keybox}
\textbf{Nominal MPC.} At every sampling time $k$:
\begin{enumerate}[label=\arabic*.]
  \item measure the state $x_k$ of the process;
  \item solve the optimization problem $\mathcal{P}_N(x_k)$ over the window $i = 0,\dots,N$,
        obtaining the optimal input sequence $\bu^*_k = \{u^*_{\ik{0}},\dots,u^*_{\ik{N-1}}\}$;
  \item apply $u_k = u^*_{\ik{0}}$ to the process and discard the rest of the sequence;
  \item wait for the next sampling time and go back to step 1.
\end{enumerate}
\end{keybox}

The first input of the optimal sequence depends on the measured state: $u^*_{\ik{0}} =
\kappa_N(x_k)$. MPC is therefore a \emph{feedback} controller, not an open-loop plan computed
once, and $\kappa_N$ is its control law. It is a law that is never written down explicitly: it is
evaluated, at every step, by solving an optimization problem. This has two consequences that run
through the whole course. First, MPC is a nonlinear controller even when the process is linear,
because $\kappa_N$ is the solution of a constrained problem. Second, the time needed to solve the
problem must be much smaller than the sampling time $T$, which is the main practical limit of
the technique.

\subsection{Why MPC, and at what price}

The strength of MPC is that the goals of the control problem are written \emph{directly} into
the optimization problem: the cost says what we want, the constraints say what is admissible.
Limits on inputs and states --- the power of a pump, the capacity of a battery, the maximum
number of vehicles a queue can hold --- are constraints, not afterthoughts; several inputs and
several outputs are handled at once; and known information about the future, such as a forecast
of the demand or the schedule of a tariff, enters the plan naturally. The formulation is also
easy to explain to whoever operates the process, because every term has a meaning.

The price is twofold. The controller needs a \emph{model} of the process, good enough for its
predictions to be useful over the window; obtaining such a model is the subject of the third
part of the course. And it needs an optimization problem to be solved at every sampling time,
which rules out, or makes expensive, the cases in which the sampling time is very short or the
problem is very large.

The problems solved by MPC belong to two families. When all the decision variables are
continuous and the cost and the constraints are convex --- a quadratic cost with linear
constraints is the typical case --- the problem is a \emph{continuous program}, solved reliably
and fast. When some decisions are logical --- start a machine or not, turn a light green or
not --- they are represented by binary variables and the problem becomes a
\emph{mixed-integer program}, much harder in general. Both families appear in the case studies.

\subsection{Centralized and decentralized control}

When the process is made of many parts --- the intersections of a city, the buildings of a
district, the vehicles of a fleet --- one can write a single optimization problem for the whole
system, which is \emph{centralized} MPC, or one problem per part, each with its own controller,
which is \emph{decentralized} MPC; when the local controllers exchange information, the scheme
is called \emph{distributed}. The centralized problem is optimal but large; the local problems
are small but have to deal with the \emph{coupling constraints} that tie the parts together.
The case studies of this part are centralized; the question comes back with the aggregation of
electric vehicles and with the control of several intersections.

\subsection{Hard and soft constraints}

A constraint is \emph{hard} when it must hold without exception: the physical limit of an
actuator, the safety limit of a tank. It is \emph{soft} when it expresses a preference that may
be violated at a cost: a reference level, a deadline for a job. Soft constraints are written by
adding a non-negative \emph{slack} variable that measures the violation and by charging it in the
cost; \cref{sec:soft} shows the construction. The distinction matters because a problem whose
hard constraints cannot all be satisfied has no solution, and a controller with no solution has
nothing to apply.

\begin{reading}
\refkey{Rawlings\_2017}, sections 1.1 and 1.2 (pp.~1--10), for the setting of the problem and for
modelling; \refkey{Cannon\_2016}, chapter 1, for an overview of predictive control. The slides on
MPC of the course website, up to slide 25, cover this chapter.
\end{reading}
